\chapter{Inhibitor Destruction by Dry Storage}

A common misconception is that seeds will germinate if given a little moisture
and a little warmth. This misconception arises because the seeds of about 50\% of all
temperate zone plants can be collected, put in an envelope on the shelf, and
germinated months later when given moisture and warmth. What is often overlooked
is that the period of dry storage on the shelf was essential to the germination. What
happens is that chemical inhibition systems are present initially and these are
destroyed by the drying. This is the commonest mechanism for preventing
germination of the seed before it is dispersed. It is convenient to call these D-70
germinators. Germinators of the D-70 type are predominant in Asteraceae (daisies),
Brassicaceae (mustards), Cam panulaceae (bellflowers), and Poaceae (grasses).

They also predominate in species of agricultural importance, because it was important
to early man to be able to dry store the seed over winter. What evolved in agriculture
was an efficient selection for this type. Most of the common garden annuals are also of
this type. The 50\% of all temperate zone plants that are not D-70 germinators tend to
be avoided and regarded as recalcitrant germinators. The use of the term immediate
germinators for these of the D-70 type is objectionable because they will not
immediately germinate unless given a prior dry storage.

The processes involved in the destruction of germination inhibitors have not
been examined from a chemical standpoint, and in fact have hardly been recognized
as chemical processes in the literature on seed germination. It is planned to determine
if oxygen is required by testing whether these processes can be conducted in sealed
envelopes and anaerobic atmospheres. It will be an exciting area of research in the
future to elucidate the nature of these chemical processes. They may prove to not be
simple destruction of a chemical compound, but rather changes in protein
conformations and changes in membrane structure. These are regarded as ultimately
chemical in nature.

The dry storage requirement and the dramatic effects of dry storage are shown
by some representative data in Table 5-1. The effect of dry storage can be overlooked
if seed is collected from seed capsules in which sufficient drying has already occurred,
particularly in species requiring only a short dry storage period. Dry storage times
have been reported to range from three weeks for barley to sixty months for Rumex
crispus (ref. 12). These timesmust depend on temperature and relative humidity.
However, in the experiments reported these variable were not precisely controlled.
Fresh seed of many 0-70 germinators show two curious behaviors. The more
important one is that if fresh seed is placed in moisture at 70, it seems to be
permanently injured and fails to give significant germination thereafter no matter what
the treatment. This was found to be true for the species in Table 5-1. The second
curious behavior was shown by the Armeria, the five Draba, and the Doronicum in
Table 5-1. Although fresh seed germinated poorly if at all at 70, it germinated fairly
well at 40 albeit slowly. It is as if the germination inhibitors that are so effective at 70
are either not soeffective at 40 or are partially destroyed at 40.

The percent germinations in Table 5-1 were calculated on the basis that the
number of seeds was the number that germinated under the optimum conditions. This
was always at least 60\% and usually 80-90\% of the normal size seed coats.

Many species of the D-70 type grow well and set abundant seed in climates
with moist summers, but few seedlings appear. The explanation is that the seed is
shed in midsummer and encounters summer moisture before the seed has had a
chance to dry. This injures the seed irreversibly. The solution is to collect the geed,
dry store it, and sow after drying. In this way colonies can be maintained.
There are some seeds that seem to germinate immediately at 70 both fresh and
dry stored. An example is Eunomia oppositifolia. It is possible that these are D-70
germinators in which the reuisite dry storage time is very short. Another species with
a related behavior is Centaurea maculosa. The seed germinates immediately at 70 on
being placed in a moist towel so the question arises as to why the seed does not
germinate in the seed capsule, particularly since the ripe seed is held in the seed
capsule for months. An explanation is given in Chapter 9 and involves a tight
receptacle that protects the seed from moisture.

\begin{table}[h!]
\centering

% \begin{tabular}{ |p{3cm}|p{1cm}|p{1cm}|p{3cm}|p{1cm}|p{1cm}| }
\begin{tabular}{|l|c|c|l|c|c|}
 \hline
 Species &
 \multicolumn{2}{|p{2cm}|}{\% Germ, 1 month at 70} &
 Species &
 \multicolumn{2}{|p{2cm}|}{\% Germ, 1 month at 70} \\
 \hline

 & fresh & dry & & fresh & dry \\

\hline
Anemone cylindrica & 4 & 98     & Draba densifolia & 0 & 100 \\
Armeria caespitosa & 15 & 90    & Draba lasiocarpa & 1 & 100 \\
Aster coloradoensis & 0 & 100   & Draba parnassica & 5 & 100 \\
Campanula carpathica & 2 & 85   & Draba sartori    & 5 & 100 \\
Doronicum caucasicum & 10 & 100 & Hesperis matronalis & 0 & 50 \\
Draba aizoon & 0 & 100          & Penstemon grandiflorus & 0 & 50 \\
Draba compacta & 2 & 100        & Townsendia parryi      & 5 & 100 \\
Draba dedeana & 0 & 100         & & & \\
\hline
\end{tabular}

\caption{Comparison of Germination of Fresh Seed Relative to Seed Dry Stored for Six Months at 70}
\label{table:fds}
\end{table}

For species of the D-70 type, two competing chemical reactions are taking place
during the dry storage. One is the destruction of germination inhibitors discussed
above. The second is steady dieing of the seeds. Ultimately all seeds die if dry
storage is prolonged far enough. The situation is best pictured as the \textit{superposition of
two plots}. Both plots have time on the horizontal coordinate. One plot proceeds
upwards from the zero-zero point as time progresses, and this represents the number
of seeds that have become ready for germination due to destruction of the germination
inhibitor. The second plot also starts from zero and is the number of seeds that have
died. The true number of germinatable seeds is the \textit{first plot minus the second plot}.
The resulting curve will go through a maximum so that there is some time when the
number of germinatable seeds maximizes. Fortunately the maximum will be a broad
one as there will be a considerable time period when the two processes approximately
compensate for each other.

In principle it is also possible to have D-40 germinators; That is the seed
requires drying to destroy inhibitors but then germinates at 40. In such species
everything is like the D-70 germinators except that the rate of germination is much
faster at 40 than at 70. However, species that germinated at 40 after dry storage
generally had induction periods of at least six weeks at 40 suggesting that further
germination inhibitors were present that had to be destroyed at 40 in moisture. Thus
the germination patterns were more complex. This type is discussed in Chapter 6.
Most garden annual and vegetables are D-70 germinators. A representative
selection were examined and these are presented here rather than in the data
summary in Chapter 20. Germinations were rapid. Celosia, marigold, portulaca,
radish, and zinnia showed emergence of the radicle in two days. Tomato, bean,
cucumber, watermelon, squash, swiss chard, corn, sweet pea, sweet William, and
candytuft showed emergence of the radicle in three days. Four to five day induction
periods were shown by alyssum, lobelia, nasturtium, and carrot. The rates were : -
approximately first order for marigold, radish, squash, sweet William, swiss chard,
watermelon, and zinnia with half lives of one to two days. The rates were
approximately zero order for candytuft, carrot, and sweet pea with rates of 25-50\%/d.
In general the germination rates were negligible at 40, but alyssum germinated at the
same rate and ind. t at 40 as at 70 and with candytuft, carrot, nasturtium, radish, sweet
william, and swiss chard the germination rate was the same but the ind??t was two to
three times longer.

The number of genera studied could have been much increased by including
studies on numerous garden annuals and Asteraceae. However, these are almost
invariably 0-70 germinators (Vernonia in Asteraceae is an exception), and were not
studied unless there was some reason for anticipating unusual behavior.

